%قبل از شروع مسئله ترجیحاً چیزی ننویسید. اما اگر خواستید بنویسید، در قالب لیستی (فقرات و شمارش) بنویسید.
%برای اضافه کردن نمره، بعدی از شروع‌مسئله، [x] را اضافه کنید و به جای x  نمره مورد نظر را قرار دهید. برای زیرمسئله‌ها هم به همین صورت است.
%\شروع‌مسئله[10]
\شروع‌مسئله
% صورت اصلی مسئله را در این‌جا بنویسید.
پیاده‌سازی درخت دودویی جست‌وجو با عناصر تکراری چالش‌هایی از نظر زمانی ایجاد می‌کند. در این مسئله می‌خواهیم به بررسی این موضوع بپردازیم. فرض کنید عمل درج ساده ما مشابه تابع $NR\_BST\_INSERT$ در صفحه ۸۵ اسلایدهاست.
%اگر مسئله چند قسمت داشت به این صورت عمل کنید:
\شروع{زیرمسائل}

\زیرمسئله
در صورتی که $n$ بار پشت سرهم بخواهیم یک داده یکسان را در درختی که در ابتدا خالی است درج کنیم، مرتبه زمانی این کار به چه صورت خواهد بود؟

در ادامه مسئله قصد داریم با چک کردن شرط یکسان بودن کلید‌ها پیش از خطوط $6$ و $14$ در الگوریتم صفحه ۸۵، به شکل‌ خاصی حالت برابری کلید‌ها را پیاده سازی کنیم. برای هر یک از حالت‌هایی که در زیر پیشنهاد می‌شود، پیچیدگی زمانی الگوریتم را برای درج $n$ عنصر یکسان تحلیل کنید.
\زیرمسئله
یک مقدار flag صفر و یکی در هر عنصر نگه داریم.آن گاه بسته به این که مقدار این flag برابر True یا False باشد، در خط 7 الگوریتم
$p \leftarrow left[p]$
یا
$p \leftarrow right[p]$
را اجرا کنیم. مقدار این flag در هر بار گذر از یک راس در هنگام درج مقداری برابر با آن راس در درخت، از True به False یا از False به True تغییر می‌کند.

\زیرمسئله
لیستی از گره‌هایی که مقادیر برابری با گره $x$ دارند را در آن نگه داریم و هر بار هنگام درج مقداری برابر با آن گره، آن مقدار را به آن لیست اضافه کنیم.

\زیرمسئله
برای این قسمت مشابه بخش ب عمل می‌کنیم با این تفاوت که دیگر flag ای در نظر نمی‌گیریم و به طور کاملا تصادفی (با احتمال برابر) به سمت چپ یا راست می‌رویم. برای این قسمت هم بدترین حالت را مشخص کنید و هر به صورت شهودی حالت میانگین را تحلیل کنید.

\پایان{زیرمسائل}
\solution{

برای قسمت اول، از آن جایی که شرط خط $14$ همواره false می‌شود، هر عنصر به راست‌ترین سمت درخت اضافه می‌شود. در نتیجه زمان اجرا به صورت
$\sum_{i=1}^{n} i \in \Theta(n^2)$
خواهد بود.

برای قسمت دوم سوال، این روش منجر می‌شود که هر کدام از زیردرختان تنها به اندازه $1$ واحد در ارتفاع با هم فاصله داشته باشند. در نتیجه ارتفاع از
$\Theta(\log(n))$
خواهد بود. در نتیجه کل زمان اجرا به صورت
$\sum_{i=1}^{n} \log(n) \in \Theta(n \log (n))$
خواهد بود.


در قسمت سوم، با فرض این که صرفا $n$ درج برابر پشت سر هم داریم، ارتفاع درخت همواره $0$ باقی می‌ماند و از آن جایی که درج ساده در لیست هم در زمان ثابت
$\Theta(1)$
امکان پذیر است، کل زمان اجرا به صورت خط و از 
$\Theta(n)$
خواهد بود.

برای قسمت چهارم، بدترین حالت این است که هر دفعه سمت راست به صورت تصادفی انتخاب شود و در این حالت مشابه بخش اول و
$\Theta(n^2)$
می‌شود. در حالت میانگین به شکل شهودی، نصف اوقات سمت راست و نصف اوقات سمت چپ را انتخاب می‌کند و در نتیجه درخت تقریبا متعادل از دو سمت گسترش پیدا می‌کند. در نتیجه مشابه بخش دوم می‌شود و در حالت میانگین به صورت
$\Theta(n \log (n))$
خواهد بود.




}



%%%%%%%%%%%%%%%%%%%%%%%%%%%%%%%%%%%%%%%%%%%%%%%%%%%%%
% برخی دستورات مفید برای تمرین‌ها. %
%%%%%%%%%%%%%%%%%%%%%%%%%%%%%%%%%%%%%%%%%%%%%%%%%%%%%

% فرمول‌های نمونه:
%$f(x) = \Theta(x)$
%$T(x, y) \leq \log(x) + 2^y + \binom{2n}{n}$
% {\tt code\_function}


% لیست موارد بدون شماره‌گذاری:
%\شروع{فقرات}
%    \فقره مورد اول 
%        \شروع{فقرات}
%            \فقره مورد اول داخلی 
%                \شروع{فقرات}
%                    \فقره اولین مورد خیلی داخلی 
%                    \فقره دومین مورد خیلی داخلی 
%                \پایان{فقرات}
%            \فقره مورد دوم داخلی 
%        \پایان{فقرات}
%    \فقره مورد دوم 
%\پایان{فقرات}

% لیست موارد شماره‌گذاری شده:
%\شروع{شمارش}
%    \فقره یک مورد 
%    \فقره یک مورد دیگر 
%    \فقره باز هم یک مورد دیگر
%\پایان{شمارش}

% برای نوشتن یک فرمول چند خطی که یک شماره بخوره. برای حالت‌های دیگر به سایت زیر مراجعه کنید:
% https://www.overleaf.com/learn/latex/Aligning_equations_with_amsmath
%\begin{align} 
%    \begin{split}
%        (x+y)^3 &= (x+y)^2(x+y) \\
%                &= (x^2+2xy+y^2)(x+y) \\
%                &= (x^3+2x^2y+xy^2) + (x^2y+2xy^2+y^3) \\
%                &= x^3+3x^2y+3xy^2+y^3
%    \end{split}                                 
%\end{align}

% نوشتن ماتریس و آرایه دوبعدی
%\begin{align}
%    A = 
%    \begin{bmatrix}
%        A_{11} & A_{21} \\
%        A_{21} & A_{22}
%    \end{bmatrix}
%\end{align}

% اضافه کردن تصویر با شرح. برای اضافه کردن تصویر بدون شرح، فقط دستور درج‌تصویر را استفاده کنید.
%\شروع{شکل}[h!]
%\تنظیم‌ازوسط
%  \درج‌تصویر[پهنا=0.5\پهنای‌متن]{img.jpg}
%\شرح{شرح تصویر}
%\برچسب{تصویر:۱}
%\پایان{شکل}
%\پایان{زیرمسائل}

%کد زیر را ببینید.
%%\lstinputlisting{source_filename.py}
%\begin{latin}
%\begin{lstlisting}
%# Import the modules
%import sys
%import random
%
%ans = True
%
%while ans:
%question = raw_input("Ask the magic 8 ball a question: (press enter to quit) ")
%
%answers = random.randint(1,8)
%
%if question == "":
%sys.exit()
%
%elif answers == 1:
%print "It is certain"
%
%elif answers == 2:
%print "Outlook good"
%
%elif answers == 3:
%print "You may rely on it"
%
%elif answers == 4:
%print "Ask again later"
%
%elif answers == 5:
%print "Concentrate and ask again"
%
%elif answers == 6:
%print "Reply hazy, try again"
%
%elif answers == 7:
%print "My reply is no"
%
%elif answers == 8:
%print "My sources say no"
%\end{lstlisting}
%\end{latin}
